%!TEX root = ../dokumentation.tex

\addchap{\langabkverz}
%nur verwendete Akronyme werden letztlich im Abkürzungsverzeichnis des Dokuments angezeigt
%Verwendung:
%		\ac{Abk.}   --> fügt die Abkürzung ein, beim ersten Aufruf wird zusätzlich automatisch die ausgeschriebene Version davor eingefügt bzw. in einer Fußnote (hierfür muss in header.tex \usepackage[printonlyused,footnote]{acronym} stehen) dargestellt
%		\acs{Abk.}   -->  fügt die Abkürzung ein
%		\acf{Abk.}   --> fügt die Abkürzung UND die Erklärung ein
%		\acl{Abk.}   --> fügt nur die Erklärung ein
%		\acp{Abk.}  --> gibt Plural aus (angefügtes 's'); das zusätzliche 'p' funktioniert auch bei obigen Befehlen
%	siehe auch: http://golatex.de/wiki/%5Cacronym
%
\begin{doublespace}
\begin{acronym}[YTMMM]
\setlength{\itemsep}{-\parsep}

\acro{CHRO}{Chief Human Resources Officer}
\acro{CIO}{Chief Information Officer}
\acro{CPI}{Cost Performance Index}
\acro{CR}{Change Request}
\acro{DSGVO}{Datenschutz-Grundverordnung}
\acro{dmTECH}{IT-Betriebsorganisation der dm-Gruppe}
\acro{EVM}{Earned Value Management}
\acro{HR}{Human Resources}
\acro{IT}{Informationstechnologie}
\acro{KPI}{Key Performance Indicator}
\acro{MOV}{Measurable Organizational Value}
\acro{NPV}{Net Present Value}
\acro{PERT}{Program Evaluation and Review Technique}
\acro{ROI}{Return on Investment}
\acro{SPI}{Schedule Performance Index}
\acro{UAT}{User Acceptance Test}
\acro{VL}{Verbundene Länder}
\acro{WBS}{Work Breakdown Structure}
\acro{WIIFM}{What's In It For Me}

\end{acronym}
\end{doublespace}
